\section*{章节综合：用不同语言解同一问题}
求复方程 $z^2-2z+2=0$。中心化后是 $(z-1)^2=-1$，所以两根为 $1+\mathrm i$ 与 $1-\mathrm i$。到这里，二次公式与复数扩展已经合在一起；接着用单位圆语言看它们，模长都是 $\sqrt2$，角度分别是正负四分之一的直角。

因此 $(1+\mathrm i)^4$ 无需展开四次：模长变成四，角度变成 $\pi$，结果为负四。用代数平方两次也得到同一结果，提供了独立检查。极坐标适合乘方，直角坐标适合加法，选择表示可以显著减少运算。

若继续取对数，主值为 $\frac12\ln2+\mathrm i\pi/4$，但所有对数还允许虚部增加 $2\pi k$。从乘方返回对数时必须保留这份角度信息，而不能把实数中的消去规则无条件搬过来。

复习本组时，给每个答案补一句“我现在使用哪一种表示、定义域和分支约定是什么”。这个习惯能把基础公式变成可迁移的推理工具。
